\section{Results}
\subsection{Sample Characteristics}

Sample characteristics are available in Table~\ref{tab:sample_characteristics}. The final sample consisted of 22 students, including 9 female and 11 male participants, while 2 participants did not specify their gender. Participants' ages ranged from 21 to 28 years ($M = 25$, $SD = 2.145$). Thirteen participants reported prior experience with Jupyter Notebook. Self-reported Python proficiency was measured on a scale from 1 to 10 and showed a relatively even distribution across different levels of expertise, allowing the dashboard to be evaluated by students with varying programming backgrounds. The participants spent an average of $17.289$ minutes testing the dashboard.
\begin{table}
  \centering
  \caption{Sample Characteristics ($N=22$)}
  \label{tab:sample_characteristics}
  \small
  \begin{tabular}{@{} l l r @{}} 
    \toprule
    \textbf{Characteristic} & \textbf{Category / Metric} & {\textbf{Value}} \\
    \midrule
    \textbf{Gender}         & Female                & 9 ($40.9\%$) \\
                            & Male                  & 11 ($50.0\%$) \\
                            & Unspecified           & 2 ($9.1\%$) \\
    \addlinespace
    \textbf{Age (Years)}    & Mean (SD)             & $25.00$ ($2.15$) \\
                            & Range                 & {21--28} \\
    \addlinespace
    \textbf{Coding Knowledge} & Mean (SD)          & $4.82$ ($2.70$) \\
                            & Range                 & {0--9} \\
    \addlinespace
    \textbf{Testing Duration (min)} & Mean (SD)   & $17.29$ ($7.68$) \\
                            & Range                 & {3--34} \\
    \bottomrule
  \end{tabular}
\end{table}

\subsection{Reliability of Scales}
The internal consistency of the multi-item scales was assessed using Cronbach's alpha~\cite{cronbach1951coefficient}. The System Usability Scale demonstrated acceptable to good reliability ($\alpha = .794$). For the Technology Acceptance Model subscales, Perceived Usefulness achieved good reliability ($\alpha = .852$), and Perceived Ease of Use showed excellent reliability ($\alpha = .941$). These reliability estimates confirm that the scales are sufficiently consistent for the purposes of this study, despite the relatively small sample size ($N = 22$). No items were removed to improve alpha, as all contributed positively to their respective scales.

\subsection{Descriptive Statistics}
Descriptive statistics for all main study variables are presented in Table~\ref{tab:descriptive_stats}.
Participants ($N = 22$) rated the perceived usability of the dashboard with a mean SUS score of $M = 65.57$ ($SD = 17.89$, range $40-90$). This value falls slightly below the widely reported average SUS benchmark of 68, indicating usability in the lower-to-average range. The Technology Acceptance Model (TAM) overall score averaged $M = 5.36$ ($SD = 0.88$, range $3.25-6.75$) on a 7-point Likert scale. The subscale Perceived Usefulness (PU) had a mean of $M = 4.93$ ($SD = 1.18$), while Perceived Ease of Use (PEOU) was higher at $M = 5.78$ ($SD = 1.06$). These values suggest moderate-to-positive perceptions of usefulness and particularly ease of use.
The subjective mental workload (NASA-TLX) yielded a mean overall score of $M = 8.15$ ($SD = 2.92$, range $1.67-12.67$) on a scale of 0 to 20. This indicates a relatively low-to-moderate perceived workload~\cite{hertzum2021reference}. 
\begin{table}
  \centering
  \caption{Descriptive Statistics of Main Study Variables ($N=22$)}
  \label{tab:descriptive_stats}
  \small
  \begin{tabular}{l r r r r}
    \toprule
    \textbf{Construct} & \textbf{M} & \textbf{SD} & \textbf{Min} & \textbf{Max} \\
    \midrule
    System Usability Scale (SUS)   & 65.57 & 17.89 & 40.00 & 90.00 \\
    TAM Overall                    &  5.36 &  0.88 &  3.25 &  6.75 \\
    \quad Perceived Usefulness (PU)  &  4.93 &  1.18 &  3.33 &  7.00 \\
    \quad Perceived Ease of Use (PEOU) &  5.78 &  1.06 &  3.17 &  7.00 \\
    NASA-TLX (Mental Workload)     &  8.15 &  2.92 &  1.67 & 12.67 \\
    \bottomrule
  \end{tabular}
\end{table}

Distributions were largely symmetric with a skewness near 0 for SUS and overall TAM, with some negative skewness in PEOU and NASA-TLX, suggesting a slight tendency toward higher ratings in those scales. Kurtosis values were acceptable overall, with no extreme deviations from normality.
To explore the relationships between usability, technology acceptance, and cognitive workload, Spearman rank-order correlations were computed (see Table~\ref{tab:correlations}). Given the small sample size ($N = 22$) and the use of Likert-scale data, non-parametric correlations were deemed appropriate. As expected, the TAM subscales Perceived Usefulness ($\rho = .816$, $p < .001$) and Perceived Ease of Use ($\rho = .706$, $p < .001$) exhibited strong, positive, and statistically significant correlations with the overall TAM score. Beyond these internal TAM associations, the analyses did not reveal any statistically significant relationships between the main constructs ($p > .05$). Specifically, perceived cognitive workload (NASA-TLX) did not significantly correlate with overall technology acceptance ($\rho = .024$, $p = .914$) or system usability ($\rho = -.099$, $p = .661$). While there was a moderate positive trend between the System Usability Scale (SUS) and Perceived Ease of Use ($\rho = .406$), this association did not reach statistical significance ($p = .061$).


\begin{table}
  \centering
  \begin{threeparttable}
    \caption{Spearman Rank Correlations ($\rho$) Between Main Study Variables ($N=22$)}
    \label{tab:correlations}
    \small
    \begin{tabular}{l r r r r r}
      \toprule
      \textbf{Variable} & \textbf{1} & \textbf{2} & \textbf{3} & \textbf{4} & \textbf{5} \\
      \midrule
      1. SUS               & ---   &       &       &       &       \\
      2. PU (TAM)          & .091  & ---   &       &       &       \\
      3. PEOU (TAM)        & .406  & .222  & ---   &       &       \\
      4. TAM Overall       & .323  & .816\tnote{*} & .706\tnote{*} & ---   &       \\
      5. NASA-TLX          & -.099 & .239  & -.255 & .024  & ---   \\
      \bottomrule
    \end{tabular}
    \begin{tablenotes}
      \footnotesize
      \item[*] $p < .001$, so these correlations are statistically significant.
    \end{tablenotes}
  \end{threeparttable}
\end{table}